\documentclass[12pt]{article}
\usepackage{amsmath, amssymb}
\usepackage{geometry}
\usepackage{hyperref}
\usepackage{physics}
\usepackage{cite}

\geometry{margin=1in}

\title{Filter Function Formalism for Qubit Noise Spectroscopy}
\author{osipra}
\date{}

\begin{document}

\maketitle

\begin{abstract}
We derive the filter function formalism for qubit noise spectroscopy from first principles. Starting from the rotating frame Hamiltonian, we connect the coherence decay of a qubit ensemble to the noise power spectral density $S(\omega)$ via the filter function $F(\omega, t)$. The central result is the relation $\chi(t) = \int_0^\infty \frac{d\omega}{2\pi} S(\omega) F(\omega, t)$, which enables reconstruction of $S(\omega)$ from measured coherence decay. Specific pulse sequence filter functions are derived in document 03.
\end{abstract}

\section{Setup: Rotating Frame Hamiltonian}

We consider a two-level system (qubit) driven on resonance. In the rotating wave approximation (RWA), the Hamiltonian is:

\begin{equation}
    \hat{H} = -\frac{\hbar \Delta}{2}\hat{\sigma}_z - \frac{\hbar \Omega_R}{2}\hat{\sigma}_x
\end{equation}

On resonance ($\Delta = 0$), with no drive ($\Omega_R = 0$) during free precession, $\hat{H} = 0$ in the ideal case. However, the qubit frequency is subject to environmental noise:

\begin{equation}
    \omega(t) = \omega_0 + \delta\omega(t)
\end{equation}

where $\delta\omega(t)$ is a stochastic fluctuation with zero mean, arising from sources such as charge noise, flux noise, and two-level system (TLS) defects in the substrate \cite{Clerk2010}.

\section{Phase Accumulation and Coherence Decay}

During free precession, the qubit accumulates a random phase:

\begin{equation}
    \varphi(t) = \int_0^t \delta\omega(t')\, dt'
\end{equation}

The qubit state on the equator of the Bloch sphere is:

\begin{equation}
    \rho = \frac{1}{2}\begin{pmatrix} 1 & e^{i\varphi} \\ e^{-i\varphi} & 1 \end{pmatrix}
\end{equation}

The off-diagonal coherence element, averaged over many noise realizations, gives the ensemble coherence:

\begin{equation}
    \langle e^{i\varphi(t)} \rangle = e^{-\chi(t)}
\end{equation}

For Gaussian noise, the cumulant expansion truncates at second order \cite{Kubo1962}:

\begin{equation}
    \langle e^{i\varphi(t)} \rangle = e^{-\frac{1}{2}\langle \varphi^2(t) \rangle}
\end{equation}

so that:

\begin{equation}
    \chi(t) = \frac{1}{2}\langle \varphi^2(t) \rangle
\end{equation}

Expanding the square of the integral:

\begin{equation}
    \langle \varphi^2(t) \rangle = \int_0^t \int_0^t \langle \delta\omega(t')\delta\omega(t'')\rangle\, dt'\, dt''
\end{equation}

\section{Wiener-Khinchin Theorem}

The autocorrelation function $C(t' - t'') = \langle \delta\omega(t')\delta\omega(t'')\rangle$ and the power spectral density $S(\omega)$ are a Fourier pair, related by the Wiener-Khinchin theorem \cite{Wiener1930}:

\begin{equation}
    C(t' - t'') = \int_{-\infty}^{\infty} \frac{d\omega}{2\pi}\, S(\omega)\, e^{i\omega(t' - t'')}
\end{equation}

Substituting into the expression for $\langle \varphi^2(t) \rangle$ and swapping the order of integration:

\begin{equation}
    \langle \varphi^2(t) \rangle = \int_{-\infty}^{\infty} \frac{d\omega}{2\pi}\, S(\omega) \left|\int_0^t e^{i\omega t'}\, dt'\right|^2
\end{equation}

\section{The Filter Function}

\subsection{General Definition}

For a general pulse sequence, the phase accumulation includes a modulation function $f(t') \in \{+1, -1\}$ that flips sign at each $\pi$ pulse \cite{Cywinski2008}:

\begin{equation}
    \varphi(t) = \int_0^t f(t')\, \delta\omega(t')\, dt'
\end{equation}

This leads to the general filter function:

\begin{equation}
    F(\omega, t) = \left|\int_0^t f(t')\, e^{i\omega t'}\, dt'\right|^2
\end{equation}

and the central relation of noise spectroscopy:

\begin{equation}
    \boxed{\chi(t) = \int_0^{\infty} \frac{d\omega}{2\pi}\, S(\omega)\, F(\omega, t)}
\end{equation}

$F(\omega, t)$ acts as a frequency-selective filter — it determines which frequencies of $S(\omega)$ contribute to the coherence decay. Different pulse sequences correspond to different modulation functions $f(t')$ and therefore different spectral windows \cite{Alvarez2011}. Explicit evaluation of $F(\omega, t)$ for specific sequences is carried out in document 03.

\subsection{Physical Interpretation}

The relation between $\chi(t)$ and $S(\omega)$ is a weighted integral — the pulse sequence acts as a spectral window into the noise environment. By designing sequences with filter functions that peak at different frequencies, one can probe $S(\omega)$ selectively across the spectrum. This is the principle of dynamical decoupling-based noise spectroscopy \cite{Bylander2011}.

The modulation function $f(t')$ encodes the entire pulse sequence:
\begin{itemize}
    \item No $\pi$ pulses: $f(t') = +1$ everywhere (Ramsey)
    \item One $\pi$ pulse at $t/2$: $f(t')$ flips once (Hahn echo)
    \item $n$ equally spaced $\pi$ pulses: $f(t')$ flips $n$ times (CPMG-$n$)
\end{itemize}

The filter function $F(\omega, t)$ is then the squared magnitude of the Fourier transform of $f(t')$, weighted by the pulse timing. This makes $F(\omega, t)$ directly computable from the pulse sequence geometry.

\begin{thebibliography}{9}

\bibitem{Cywinski2008}
L. Cywi\'{n}ski, R. M. Lutchyn, C. P. Nave, and S. Das Sarma,
\textit{How to enhance dephasing time in superconducting qubits},
Phys. Rev. B \textbf{77}, 174509 (2008).
\href{https://arxiv.org/abs/0807.4926}{arXiv:0807.4926}

\bibitem{Alvarez2011}
G. A. \'{A}lvarez and D. Suter,
\textit{Measuring the Spectrum of Colored Noise by Dynamical Decoupling},
Phys. Rev. Lett. \textbf{107}, 230501 (2011).

\bibitem{Bylander2011}
J. Bylander et al.,
\textit{Noise spectroscopy through dynamical decoupling with a superconducting flux qubit},
Nature Physics \textbf{7}, 565 (2011).

\bibitem{Clerk2010}
A. A. Clerk, M. H. Devoret, S. M. Girvin, F. Marquardt, and R. J. Schoelkopf,
\textit{Introduction to quantum noise, measurement, and amplification},
Rev. Mod. Phys. \textbf{82}, 1155 (2010).

\bibitem{Kubo1962}
R. Kubo,
\textit{Generalized Cumulant Expansion Method},
J. Phys. Soc. Japan \textbf{17}, 1100 (1962).

\bibitem{Wiener1930}
N. Wiener,
\textit{Generalized harmonic analysis},
Acta Math. \textbf{55}, 117 (1930).

\end{thebibliography}

\end{document}